\documentclass[a4paper,12pt,twoside, openright]{memoir}
\usepackage[left=4cm, right=3cm]{geometry}

% Castellano
\usepackage[spanish,es-tabla]{babel}
\selectlanguage{spanish}
\usepackage[utf8]{inputenc}
\usepackage[T1]{fontenc}
\usepackage{lmodern} % Scalable font
\usepackage{microtype}
\usepackage{placeins}

\RequirePackage{booktabs}
\RequirePackage[table]{xcolor}
\RequirePackage{xtab}
\RequirePackage{multirow}

% Links
\PassOptionsToPackage{hyphens}{url}\usepackage[colorlinks]{hyperref}
\hypersetup{
	allcolors = {red}
}

% Ecuaciones
\usepackage{amsmath}

% Rutas de fichero / paquete
\newcommand{\ruta}[1]{{\sffamily #1}}

% Párrafos
\nonzeroparskip

% Huérfanas y viudas
\widowpenalty100000
\clubpenalty100000

% Imágenes

% Comando para insertar una imagen en un lugar concreto.
% Los parámetros son:
% 1 --> Ruta absoluta/relativa de la figura
% 2 --> Texto a pie de figura
% 3 --> Tamaño en tanto por uno relativo al ancho de página
\usepackage{graphicx}
\newcommand{\imagen}[3]{
	\begin{figure}[!h]
		\centering
		\includegraphics[width=#3\textwidth]{#1}
		\caption{#2}\label{fig:#1}
	\end{figure}
	\FloatBarrier
}

% Comando para insertar una imagen sin posición.
% Los parámetros son:
% 1 --> Ruta absoluta/relativa de la figura
% 2 --> Texto a pie de figura
% 3 --> Tamaño en tanto por uno relativo al ancho de página
\newcommand{\imagenflotante}[3]{
	\begin{figure}
		\centering
		\includegraphics[width=#3\textwidth]{#1}
		\caption{#2}\label{fig:#1}
	\end{figure}
}

% El comando \figura nos permite insertar figuras comodamente, y utilizando
% siempre el mismo formato. Los parametros son:
% 1 --> Porcentaje del ancho de página que ocupará la figura (de 0 a 1)
% 2 --> Fichero de la imagen
% 3 --> Texto a pie de imagen
% 4 --> Etiqueta (label) para referencias
% 5 --> Opciones que queramos pasarle al \includegraphics
% 6 --> Opciones de posicionamiento a pasarle a \begin{figure}
\newcommand{\figuraConPosicion}[6]{%
  \setlength{\anchoFloat}{#1\textwidth}%
  \addtolength{\anchoFloat}{-4\fboxsep}%
  \setlength{\anchoFigura}{\anchoFloat}%
  \begin{figure}[#6]
    \begin{center}%
      \Ovalbox{%
        \begin{minipage}{\anchoFloat}%
          \begin{center}%
            \includegraphics[width=\anchoFigura,#5]{#2}%
            \caption{#3}%
            \label{#4}%
          \end{center}%
        \end{minipage}
      }%
    \end{center}%
  \end{figure}%
}

%
% Comando para incluir imágenes en formato apaisado (sin marco).
\newcommand{\figuraApaisadaSinMarco}[5]{%
  \begin{figure}%
    \begin{center}%
    \includegraphics[angle=90,height=#1\textheight,#5]{#2}%
    \caption{#3}%
    \label{#4}%
    \end{center}%
  \end{figure}%
}
% Para las tablas
\newcommand{\otoprule}{\midrule [\heavyrulewidth]}
%
% Nuevo comando para tablas pequeñas (menos de una página).
\newcommand{\tablaSmall}[5]{%
 \begin{table}
  \begin{center}
   \rowcolors {2}{gray!35}{}
   \begin{tabular}{#2}
    \toprule
    #4
    \otoprule
    #5
    \bottomrule
   \end{tabular}
   \caption{#1}
   \label{tabla:#3}
  \end{center}
 \end{table}
}

%
% Nuevo comando para tablas pequeñas (menos de una página).
\newcommand{\tablaSmallSinColores}[5]{%
 \begin{table}[H]
  \begin{center}
   \begin{tabular}{#2}
    \toprule
    #4
    \otoprule
    #5
    \bottomrule
   \end{tabular}
   \caption{#1}
   \label{tabla:#3}
  \end{center}
 \end{table}
}

\newcommand{\tablaApaisadaSmall}[5]{%
\begin{landscape}
  \begin{table}
   \begin{center}
    \rowcolors {2}{gray!35}{}
    \begin{tabular}{#2}
     \toprule
     #4
     \otoprule
     #5
     \bottomrule
    \end{tabular}
    \caption{#1}
    \label{tabla:#3}
   \end{center}
  \end{table}
\end{landscape}
}

%
% Nuevo comando para tablas grandes con cabecera y filas alternas coloreadas en gris.
\newcommand{\tabla}[6]{%
  \begin{center}
    \tablefirsthead{
      \toprule
      #5
      \otoprule
    }
    \tablehead{
      \multicolumn{#3}{l}{\small\sl continúa desde la página anterior}\\
      \toprule
      #5
      \otoprule
    }
    \tabletail{
      \hline
      \multicolumn{#3}{r}{\small\sl continúa en la página siguiente}\\
    }
    \tablelasttail{
      \hline
    }
    \bottomcaption{#1}
    \rowcolors {2}{gray!35}{}
    \begin{xtabular}{#2}
      #6
      \bottomrule
    \end{xtabular}
    \label{tabla:#4}
  \end{center}
}

%
% Nuevo comando para tablas grandes con cabecera.
\newcommand{\tablaSinColores}[6]{%
  \begin{center}
    \tablefirsthead{
      \toprule
      #5
      \otoprule
    }
    \tablehead{
      \multicolumn{#3}{l}{\small\sl continúa desde la página anterior}\\
      \toprule
      #5
      \otoprule
    }
    \tabletail{
      \hline
      \multicolumn{#3}{r}{\small\sl continúa en la página siguiente}\\
    }
    \tablelasttail{
      \hline
    }
    \bottomcaption{#1}
    \begin{xtabular}{#2}
      #6
      \bottomrule
    \end{xtabular}
    \label{tabla:#4}
  \end{center}
}

%
% Nuevo comando para tablas grandes sin cabecera.
\newcommand{\tablaSinCabecera}[5]{%
  \begin{center}
    \tablefirsthead{
      \toprule
    }
    \tablehead{
      \multicolumn{#3}{l}{\small\sl continúa desde la página anterior}\\
      \hline
    }
    \tabletail{
      \hline
      \multicolumn{#3}{r}{\small\sl continúa en la página siguiente}\\
    }
    \tablelasttail{
      \hline
    }
    \bottomcaption{#1}
  \begin{xtabular}{#2}
    #5
   \bottomrule
  \end{xtabular}
  \label{tabla:#4}
  \end{center}
}



\definecolor{cgoLight}{HTML}{EEEEEE}
\definecolor{cgoExtralight}{HTML}{FFFFFF}

%
% Nuevo comando para tablas grandes sin cabecera.
\newcommand{\tablaSinCabeceraConBandas}[5]{%
  \begin{center}
    \tablefirsthead{
      \toprule
    }
    \tablehead{
      \multicolumn{#3}{l}{\small\sl continúa desde la página anterior}\\
      \hline
    }
    \tabletail{
      \hline
      \multicolumn{#3}{r}{\small\sl continúa en la página siguiente}\\
    }
    \tablelasttail{
      \hline
    }
    \bottomcaption{#1}
    \rowcolors[]{1}{cgoExtralight}{cgoLight}

  \begin{xtabular}{#2}
    #5
   \bottomrule
  \end{xtabular}
  \label{tabla:#4}
  \end{center}
}



\graphicspath{ {./img/} }

% Capítulos
\chapterstyle{bianchi}
\newcommand{\capitulo}[2]{
	\setcounter{chapter}{#1}
	\setcounter{section}{0}
	\setcounter{figure}{0}
	\setcounter{table}{0}
	\chapter*{#2}
	\addcontentsline{toc}{chapter}{#2}
	\markboth{#2}{#2}
}

% Apéndices
\renewcommand{\appendixname}{Apéndice}
\renewcommand*\cftappendixname{\appendixname}

\newcommand{\apendice}[1]{
	%\renewcommand{\thechapter}{A}
	\chapter{#1}
}

\renewcommand*\cftappendixname{\appendixname\ }

% Formato de portada
\makeatletter
\usepackage{xcolor}
\newcommand{\tutor}[1]{\def\@tutor{#1}}
\newcommand{\course}[1]{\def\@course{#1}}
\definecolor{cpardoBox}{HTML}{E6E6FF}
\def\maketitle{
  \null
  \thispagestyle{empty}
  % Cabecera ----------------
\noindent\includegraphics[width=\textwidth]{cabecera}\vspace{1cm}%
  \vfill
  % Título proyecto y escudo informática ----------------
  \colorbox{cpardoBox}{%
    \begin{minipage}{.8\textwidth}
      \vspace{.5cm}\Large
      \begin{center}
      \textbf{TFG del Grado en Ingeniería Informática}\vspace{.6cm}\\
      \textbf{\LARGE\@title{}}
      \end{center}
      \vspace{.2cm}
    \end{minipage}

  }%
  \hfill\begin{minipage}{.20\textwidth}
    \includegraphics[width=\textwidth]{escudoInfor}
  \end{minipage}
  \vfill
  % Datos de alumno, curso y tutores ------------------
  \begin{center}%
  {%
    \noindent\LARGE
    Presentado por \@author{}\\ 
    en Universidad de Burgos --- \@date{}\\
    Tutor: \@tutor{}\\
  }%
  \end{center}%
  \null
  \cleardoublepage
  }
\makeatother

\newcommand{\nombre}{Estela Victoria Ballester Delgado} %%% cambio de comando

% Datos de portada
\title{TeachMePlay}
\author{\nombre}
\tutor{José Manuel Galán Ordax y Virginia Ahedo García}
\date{\today}

\begin{document}

\maketitle


\newpage\null\thispagestyle{empty}\newpage


%%%%%%%%%%%%%%%%%%%%%%%%%%%%%%%%%%%%%%%%%%%%%%%%%%%%%%%%%%%%%%%%%%%%%%%%%%%%%%%%%%%%%%%%
\thispagestyle{empty}


\noindent\includegraphics[width=\textwidth]{cabecera}\vspace{1cm}

\noindent D. José Manuel Galán Ordax y Dña. Virginia Ahedo García, profesores del
departamento de Ingeniería de Organización, área de Organización de Empresas.

\noindent Exponen:

\noindent Que el alumno Dña. \nombre, con DNI 71565429D, ha realizado el Trabajo final de Grado en Ingeniería Informática titulado TeachMePlay. 

\noindent Y que dicho trabajo ha sido realizado por el alumno bajo la dirección del que suscribe, en virtud de lo cual se autoriza su presentación y defensa.

\begin{center} %\large
En Burgos, {\large \today}
\end{center}

\vfill\vfill\vfill

% Author and supervisor
\begin{minipage}{0.45\textwidth}
\begin{flushleft} %\large
Vº. Bº. del Tutor:\\[2cm]
D. nombre tutor
\end{flushleft}
\end{minipage}
\hfill
\begin{minipage}{0.45\textwidth}
\begin{flushleft} %\large
Vº. Bº. del co-tutor:\\[2cm]
D. nombre co-tutor
\end{flushleft}
\end{minipage}
\hfill

\vfill

% para casos con solo un tutor comentar lo anterior
% y descomentar lo siguiente
%Vº. Bº. del Tutor:\\[2cm]
%D. nombre tutor


\newpage\null\thispagestyle{empty}\newpage




\frontmatter

% Abstract en castellano
\renewcommand*\abstractname{Resumen}
\begin{abstract}
Este proyecto es una aplicación web que consiste en un repositorio de juegos docentes. Permite a los usuarios, según su rol de acceso, buscar y filtrar juegos, interactuar mediante valoraciones, y también añadir y modificar nuevos juegos.

La aplicación web es accesible a través del enlace: \url{https://teachmeplay.herokuapp.com/}
\end{abstract}

\renewcommand*\abstractname{Descriptores}
\begin{abstract}
Gamificación, aplicación web, Flask, Python, PostgreSQL, HTML, CSS, Heroku, SCRUM, Overleaf, GitHub\ldots
\end{abstract}

\clearpage

% Abstract en inglés
\renewcommand*\abstractname{Abstract}
\begin{abstract}
This project is a web application that consists of a repository of educational games. It allows users, according to their access role, to search and filter games, interact through ratings, and also add and modify new games.

The web application is accessible through the link: \url{https://teachmeplay.herokuapp.com/}
\end{abstract}

\renewcommand*\abstractname{Keywords}
\begin{abstract}
Gamification, web application, Flask, Python, PostgreSQL, HTML, CSS, Heroku, SCRUM, Overleaf, GitHub\ldots
\end{abstract}

\clearpage

% Indices
\tableofcontents

\clearpage

\listoffigures

\clearpage

\listoftables
\clearpage

\mainmatter
\capitulo{1}{Introducción}

La gran cantidad de horas lectivas que se invierten en las formaciones académicas puede ser un obstáculo en el aprendizaje de los alumnos, ya que disminuye el grado de concentración durante las clases. Por ello, en los últimos años han surgido dos nuevos conceptos como posible solución a este problema: gamificación y serious games.

\cite{gamificacion} La gamificación es un sistema de enseñanza que busca el aprendizaje a través de juegos. El principal objetivo que pretende conseguir es obtener la motivación de los alumnos a través de sistemas de puntuación que les proporcionen puntos de recompensa al cumplir deferentes propósitos. De esta forma, se emplean los medios y mecanismos de juego para aumentar tanto el conocimiento académico, como sus resultados. 

Por ello, las principales ventajas que proporciona son la motivación, el aumento del rendimiento y de la colaboración, la utilización y familiarización de las TICs, junto con el progreso personal.

Los avances tecnológicos y la aparición del e-learning están permitiendo el aumento de la creación y disposición de los cada vez más conocidos como juegos docentes. Actualmente, cada vez se encuentran disponibles más herramientas educativas. Entre las mas conocidas se incluyen: Kahoot, Edmodo, Socrative, Quizizz, entre otras.

\cite{serious} Los serious games son aquellos juegos que se emplean principalmente en diferentes ámbitos de la enseñanza con el objetivo de que los estudiantes aprendan de una forma más interactiva y divertida. 

\cite{gamyserious} La gamificación y los serious games se diferencian fundamentalmente en que la gamificación no se considera un juego, sino que emplea mecanismos de juego con el fin de obtener puntos de recompensa al cumplir objetivos. En cambio, los serious games son juegos completos que pretenden el aprendizaje y formación en distintas áreas.

Este proyecto tiene como objetivo la creación de una aplicación web que permita tanto la visualización como la interacción con el conjunto de juegos docentes relevantes para la docencia de distintas materias. De esta forma, al disponer de los juegos cargados en la base de datos se facilitará su gestión y organización entre los docentes, además gracias a la interacción mediante sistemas de puntuación posibilitará a los usuarios búsquedas adaptadas a sus intereses.

\section{Estructura de la memoria}
En la memoria se encuentran los siguientes apartados:
\begin{itemize}
    \item \textbf{Introducción:} presentación del tema del proyecto que se va a desarrollar. Se incluye la estructura de la memoria y la estructura de los anexos.
    \item \textbf{Objetivos del proyecto:} definición de los objetivos generales, técnicos y personales a alcanzar durante el desarrollo del proyecto.
    \item \textbf{Conceptos teóricos:} explicación de los conceptos teóricos fundamentales para la adecuada comprensión del proyecto.
    \item \textbf{Técnicas y herramientas:} exposición de las técnicas y herramientas que se van a emplear para el desarrollo y cumplimiento de los objetivos establecidos del proyecto.
    \item \textbf{Aspectos relevantes del desarrollo:} muestra de las etapas más notables durante la elaboración del proyecto. 
    \item \textbf{Trabajos relacionados:} exposición de distintos trabajos y proyectos ya finalizados en el campo del proyecto en desarrollo.
    \item \textbf{Conclusiones y líneas de trabajo futuras:} presentación de las conclusiones obtenidas al final del desarrollo, junto con posibles mejoras que se puedan aplicar al resultado final.
\end{itemize}
\newpage

\section{Estructura de los anexos}
En el anexo se encuentran los siguientes apartados:
\begin{itemize}
    \item \textbf{Plan de proyecto de software:} presentación de la planificación temporal del proyecto junto con su análisis económico y legal.
    \item \textbf{Especificación de requisitos:} descripción de los requisitos establecidos que el proyecto debe alcanzar.
    \item \textbf{Especificación de diseño:} muestra el diseño de datos, diseño procedimental, diseño arquitectónico e interfaces.
    \item \textbf{Manual del programador:} guía que incluye los aspectos más importantes respecto al código fuente del proyecto.
    \item \textbf{Manual del usuario:} guía que explica al usuario el uso adecuado de la aplicación.
\end{itemize}

\section{Materiales adjuntos}
Además de la documentación que incluye la memoria y los anexo, se adjuntan los siguientes materiales:
\begin{itemize}
    \item Repositorio en GitHub.
    \item Código del proyecto.
    \item Proyecto desplegado.  
    \item Máquina virtual. 
\end{itemize}
\capitulo{2}{Objetivos del proyecto}

En este apartado se explican los objetivos que se pretenden alcanzar durante el desarrollo del proyecto.

Estos objetivos se clasifican en objetivos generales y objetivos técnicos, siendo estos últimos los necesarios para el cumplimiento de los objetivos generales.

\section {Objetivos generales}
Los objetivos generales establecidos son los siguientes:
\begin{itemize}
    \item Crear una aplicación web que ofrezca una amplia variedad de juegos educativos para su uso en el aula, incluyendo juegos para diversas materias y niveles educativos.
    \item Desarrollar una base de datos bien estructurada y eficiente que almacene toda la información relevante sobre los juegos docentes, incluyendo detalles sobre los juegos, las cuentas de usuario, las valoraciones de los juegos y las solicitudes de rol.
    \item Proporcionar a los usuarios diferentes roles y permisos según sus necesidades y objetivos, incluyendo roles de usuario, profesor y administrador, cada uno con diferentes niveles de acceso y funcionalidades.
    \item Permitir que los usuarios interactúen con la aplicación y con los juegos mediante puntuaciones y comentarios.
    \item Fomentar el aprendizaje y la educación a través de juegos educativos, proporcionando una herramienta efectiva para el aprendizaje.
\end{itemize}

\section {Objetivos técnicos}
Los objetivos técnicos establecidos son los siguientes:
\begin{itemize}
    \item Garantizar la disponibilidad y accesibilidad de la aplicación en cualquier momento siempre que se tenga conexión a la red.
    \item Crear una interfaz visualmente atractiva que facilite su uso a los usuarios.
    \item Crear una aplicación interactiva que permita a los usuarios realizar un feedback.
    \item Aplicar la metodología ágil Scrum para la organización del desarrollo del proyecto.
    \item Emplear el gestor de referencias bibliográficas Zotero.
    \item Usar la herramienta Zube para la planificación, seguimiento y gestión de las tareas y objetivos del proyecto
\end{itemize}
	
\section{Objetivos personales}
Los objetivos personales establecidos son los siguientes:
\begin{itemize}
    \item Aprender a utilizar el framework Flask en Python para crear aplicaciones web desde cero.
    \item Mejorar mi conocimiento y habilidades en múltiples lenguajes de programación como Python, HTML, CSS y SQL.
    \item Desplegar la aplicación en línea para que pueda ser accesible desde cualquier lugar.
\end{itemize}
\capitulo{3}{Conceptos teóricos}

En esta sección, se presentan los conceptos teóricos clave que sustentan el desarrollo del proyecto. Estos conceptos son fundamentales para comprender el contexto en el que se desarrolla el proyecto y su importancia en el ámbito en el que se enmarca.

\section{Gamificación}
\cite{gamificacion1} La gamificación es un método de aprendizaje que emplea la mecánica de los juegos en el ámbito educativo y profesional, con el objetivo de aumentar y mejorar los resultados mediante el uso de entornos digitales y educativos.

Gracias a su carácter lúdico, la gamificación facilita la obtención de conocimientos, ya que ofrece una experiencia positiva al usuario y su objetivo principal es la motivación de los estudiantes. Para ello, se emplean técnicas mecánicas y dinámicas de los juegos.

\cite{gamificacion2} La técnicas mecánicas son aquellas que buscan recompensar al usuario según el cumplimiento de los requisitos del juego. Alguno de los elementos de compensación empleados son:
\begin{itemize}
    \item Clasificaciones: clasificar a los estudiantes en función de sistemas de puntos o logros.
    \item Acumulación de puntos: asignar valores cuantitativos en función de tareas logradas.
    \item Misiones o retos: superación de objetivos establecidos.
    \item Escalado de niveles: superación de niveles según los objetivos logrados.
\end{itemize}

Mientras que las técnicas dinámicas son aquellas que buscan la motivación del usuario para que continúe y complete sus objetivos. Algunas de las técnicas dinámicas empleadas son:
\begin{itemize}
    \item Estatus: establecer niveles jerárquicos.
    \item Recompensa: obtener beneficios.
    \item Logro: superación personal.
\end{itemize}

De esta forma la integración de la gamificación en entornos educativos y profesionales, hace que se creen experiencias más atractivas e interactivas para los alumnos, por lo que aumenta su compromiso y rendimiento.

Así, la gamificación se conoce como un método innovador y eficaz para mejorar la motivación y el rendimiento en el ámbito educativo.

\section{Serious games}
\cite{seriousgames} Los serious games son juegos diseñados para proporcionar formación en diversos ámbitos, como la educación, la ciencia y la ingeniería. Este concepto surge de la investigación de las distintas modalidades en las que los juegos pueden incluir métodos de aprendizaje y enseñanza sin comprometer la diversión.

Las estrategias lúdicas empleadas por los serious games permiten que los alumnos se involucren más en los contenidos educativos y disfruten de una experiencia de aprendizaje más atractiva y motivadora.

En España, existen diversas empresas e iniciativas que utilizan los serious games para la formación educativa, tales como Humantiks, Gamelearn®, Binnakle, Netlanguages y Chiara.


\section{E-learning}
\cite{e-learning} El e-learning es un concepto que se refiere a las actividades formativas que se realizan a través de dispositivos conectados a la red, lo que lo convierte en una forma de aprendizaje conocida como formación en línea o virtual. Esta modalidad facilita la educación online y permite a los estudiantes la formación a distancia.

Las ventajas y características más destacadas del e-learning son la flexibilidad que proporciona a los estudiantes y la riqueza de aprendizaje que ofrece la amplia variedad de recursos educativos que se encuentran en internet. Estos recursos abarcan todo tipo de modalidades, desde vídeos y juegos hasta herramientas interactivas y simulaciones.

\cite{e-learning2} La evolución del e-learning ha llevado al desarrollo de plataformas avanzadas, como Moodle, un software que ha estandarizado la creación de distintas plataformas de e-learning y ha permitido a los profesores crear entornos virtuales para la interacción y compartición de recursos docentes.

El e-learning ha impulsado la creación de nuevas estrategias, como el aprendizaje basado en la gamificación, en proyectos o la personalización del aprendizaje. Estrategias que permiten a los estudiantes aprender de manera más lúdica y participativa, y a los profesores adaptarse a las necesidades de los estudiantes.




\capitulo{4}{Técnicas y herramientas}

Esta parte de la memoria tiene como objetivo presentar las técnicas metodológicas y las herramientas de desarrollo que se han utilizado para llevar a cabo el proyecto. Si se han estudiado diferentes alternativas de metodologías, herramientas, bibliotecas se puede hacer un resumen de los aspectos más destacados de cada alternativa, incluyendo comparativas entre las distintas opciones y una justificación de las elecciones realizadas. 
No se pretende que este apartado se convierta en un capítulo de un libro dedicado a cada una de las alternativas, sino comentar los aspectos más destacados de cada opción, con un repaso somero a los fundamentos esenciales y referencias bibliográficas para que el lector pueda ampliar su conocimiento sobre el tema.
-----------------------------------------------------------------------------

En este apartado se muestran las distintas técnicas y herramientas empleadas para el desarrollo del proyecto.

Entre estas técnicas y herramientas se distinguen las utilizadas para la gestión del proyecto, la documentación y los entornos de desarrollo.

\section{Gestión del proyecto}
\subsection{Scrum}
\cite{Scrum1_bib} Es un marco de trabajo colaborativo que permite trabajar en equipo el desarrollo de un proyecto sistematizado en entregas organizadas parcialmente del producto final, en función de la utilidad que ofrecen al destinatario del proyecto. 

\cite{Scrum2_bib} Esta técnica de metodología ágil permite el avance de los requisitos establecidos del proyecto en planes temporales de duración corta y fija, de manera que a los requisitos se les establece una prioridad según su importancia. De esta forma, permite al equipo organizarse y sincronizarse para alcanzar sus objetivos y mostrar al cliente los resultados obtenidos al finalizar cada sprint.
\imagen{scrum}{Metodología ágil: Scrum}{0.7}

\subsection{GitHub}
\cite{GitHub1_bib} Es una plataforma que permite la creación de repositorios de código con el fin de almacenarlos en la nube mediante un sistema de control de versiones conocido como Git. Permite la gestión de proyectos de forma cooperativa en tiempo real.

\cite{GitHub2_bib} Facilita el trabajo cooperativo proporcionando distintas herramientas para la gestión de las versiones de las páginas, un sistema para el control de problemas, herramientas para la comprobación de código y un visualizador de las ramas.
\imagen{GitHub-Logo}{GitHub-logo}{0.1}

\subsection{GitHub Desktop}
\cite{GitHubDesktop_bib} Es una aplicación que posibilita la interacción con GitHub proporcionando una interfaz gráfica de usuario en sustitución a la línea de comandos o buscadores web.

Se trata de la versión de escritorio de GitHub y permite la subida, extracción y clonación de repositorios empleando mecanismos colaborativos.
\imagen{GitHub Dektop-logo}{GitHub Desktop-logo}{0.1}

\subsection{ZenHub}
\cite{ZenHub_bib} Es una plataforma de gestión de proyectos integrada en GitHub que permite la conexión de repositorios y la organización de los proyectos mediante tableros de trabajo que facilitan la planificación de sprints, la creación de tareas, la estimación de story points y la generación de gráficos e informes estadísticos.
\imagen{zenhub-logo}{ZenHub-logo}{0.1}

\subsection{Zube}
\cite{Zube_bib} Es un sistema de gestión de proyectos gratuito integrado en GitHub que facilita la asignación y gestión de tareas, permitiendo especificar la prioridad de cada una y promoviendo la colaboración entre los miembros del equipo. Además, su funcionalidad de actualización en tiempo real permite que los cambios realizados se reflejen instantáneamente en el proyecto. Por otra parte, esta plataforma se enfoca en la detección y corrección de errores, contribuyendo a garantizar la calidad y eficiencia del proyecto en todo momento.
\imagen{zube-logo}{Zube-logo}{0.1}

\section{Herramientas de documentación}
\subsection{LaTeX}
\cite{LaTeX_bib} Es un sistema de composición de textos basado en TeX. Su función consiste en la creación
de distintos documentos como libros o artículos técnicos y científicos, proporcionando una alta calidad tipográfica. 

Este procesador de textos se basa en instrucciones, por lo que la principal diferencia que ofrece respecto a otros procesadores de texto es la facilidad de incluir el contenido sin necesidad de manejar los detalles del formato permitiendo estructurar de forma muy sencilla el documento, gracias además a sus capacidades gráficas.
\imagen{latex-logo}{LaTeX-logo}{0.1}

\subsection{Overleaf}
\cite{Overleaf_bib} Es una herramienta en línea para la redacción colaborativa y su publicación. Gracias al editor LaTeX permite la cooperación en tiempo real y la compilación de la salida de forma automática en segundo plano durante la redacción de los textos.
\imagen{overleaf-logo}{Overleaf-logo}{0.1}

\subsection{Zotero}
\cite{Zotero_bib} Es un gestor de referencias bibliográficas gratuito y de código abierto que permite la recopilación y organización de distintos recursos, incluyendo su citación, sincronización y colaboración. 
\imagen{zotero-logo}{Zotero-logo}{0.1}

\section{Heramientas de diseño}
\subsection{Figma}
\cite{Figma1_bib} Es una herramienta de prototipado principalmente web que permite el diseño de las interfaces de usuario y el diseño de experiencia de las aplicaciones, posibilitando la colaboración en tiempo real.

\cite{Figma2_bib} La principal ventaja de esta herramienta es la posibilidad que ofrece de usar funcionalidades y librerías de componentes creadas y compartidas por la comunidad, lo cual facilita y agiliza el diseño de las interfaces.
\imagen{figma-logo}{Figma-logo}{0.1}

\section{Entornos de desarrollo}
 \subsection{Python}
\cite{Python1_bib} Es un lenguaje de programación de alto nivel empleado mayoritariamente para el desarrollo de software. \cite{Python2_bib} Al ser un lenguaje interpretado no es necesaria la compilación para la ejecución de las aplicaciones ya que dispone del programa conocido como interpretador, y por lo tanto no requiere la traducción a lenguaje máquina.
\imagen{Python-logo}{Python-logo}{0.1}

\subsection{Flask}
\cite{Flask_bib} Es un micro framework web que permite la creación de aplicaciones web con Python. Se basa en el Modelo-Vista-Controlador (MVC) facilitando el desarrollo de las aplicaciones, y diferenciando el modelo de datos, la vista y el controlador.
Esta herramienta proporciona utilidades para el desarrollo de aplicaciones web dinámicas permitiendo su comunicación con bases de datos, la autenticación de usuarios y el desarrollo de formularios, entre otros.
\imagen{flask-logo}{Flask-logo}{0.1}

\subsection{PostgreSQL} 
\cite{PostgreSQL_bib} Es un sistema de gestión de bases de datos relacionales de código abierto y gratuito.
PostgreSQL es conocido por su escalabilidad y la facilidad de administrar grandes cantidades de datos, siendo compatible con herramientas y tecnologías empleados en el proyecto, incluyendo lenguajes de programación como Python y frameworks web, como Flask.
\imagen{postgreSQL-logo}{PostgreSQL-logo}{0.1}

\section{Librerías}


\capitulo{5}{Aspectos relevantes del desarrollo del proyecto}

En este apartado se pretende comentar las fases más interesantes y significativas del desarrollo del proyecto, incluyendo sus inicios, la metodología empleada, la creación, el despliegue de la aplicación y la resolución de los problemas que han surgido durante el desarrollo. 


\section{Inicio del proyecto}
El tema del proyecto fue propuesto por José Ignacio Santos Martin y José Manel Galán Ordax, ya que tuvimos una primera reunión en la que me ofrecieron la posibilidad de realizar este trabajo de fin de grado, que sería tutelado por José Manuel Galán y Virginia Ahedo.

La primera propuesta fue la creación de una base de datos que se desplegase en GitHub Pages mediante JavaScript, que permitiese la subida de contenidos, su búsqueda y consulta, y la recolección de opiniones de usuarios a través de un sistema de rating a semejanza de Amazon o similares.

\section {Primeros pasos}
En la siguiente reunión que tuve con mis tutores me propusieron los primeros pasos a realizar. 

Inicialmente me plantearon gestionar el proyecto mediante la metodología ágil. Para ello, se creó un repositorio en GitHub para llevar el seguimiento del proyecto a través de sprints de dos semanas de duración.

Durante este período se empezó con la búsqueda, selección y familiarización de las distintas herramientas que se iban a utilizar para el desarrollo del proyecto. 

Algunas de las recomendaciones iniciales que me propusieron mis tutores fueron emplear un gestor de referencias como Zotero, y un gestor de base de datos como MySQL, MongoDB, MariaDB o PostgreSQL.

Respecto a la propuesta inicial de desarrollar el proyecto usando JavaScript, me decanté por usar como lenguaje de programación Python, ya que era el lenguaje aprendido durante los últimos años del grado, y por tanto iba a facilitar el desarrollo del proyecto puesto que JavaScript era un lenguaje de programación que no había usado por el momento.

Los frameworks que mis tutores me recomendaron fueros varios, como Flask o Django. Al ser Python el lenguaje de programación elegido para el desarrollo decidí usar Flask como framework. Aunque son parecidos se considera que es más conveniente Flask al facilitar el desarrollo de aplicaciones web en proyectos de menor escala como este.

\section{Metodologías}
Para la gestión del proyecto se ha empleado la metodología conocida como SCRUM, ya que fue la aprendida en el curso anterior en la asignatura de gestión de proyectos. 

Esta metodología permite la gestión del proyecto mediante el establecimiento de tareas trabajando colaborativamente. Por ello, no se ha aplicado estrictamente esta metodología ya que el proyecto no se desarrolla en equipo, pero sí que se han seguido la mayoría de prácticas que incluye la metodología ágil.

Para la división del desarrollo del proyecto se han usado los sprints que han permitido la organización de las issues en períodos de dos semanas. En estos sprints se marcaban las tareas a realizar para alcanzar los objetivos establecidos, junto con sus story points correspondientes. Para ello, se ha utilizado la herramienta ZenHub, que gracias al tablero canvas que proporciona, se pueden gestionar los sprints, tareas, etiquetas y story points, entre otros. 

Una vez finalizado el sprint se realizaba una reunión con los tutores en la que se realizaban una serie de tareas:
\begin{itemize}
\item \textbf{Revisión del sprint:} en esta tarea se revisaba si se habían cumplido todas las tareas asignadas al sprint, comprobando los avances obtenidos y los problemas que habían surgido durante este.
\item \textbf{Planificación del siguiente sprint:} una vez revisado el sprint finalizado se planificaban los siguientes objetivos a alcanzar en el siguiente sprint, asignando las tareas y estimando sus correspondientes story points.
\end{itemize}

\section{Formación}
Para llevar a cabo el desarrollo del proyecto se ha realizado la investigación, y por tanto la familiarización de distintos conceptos y herramientas. 
\subsection{Flask}

\subsection{HTML Y CSS}

\subsection{Bootstrap}

\section{Desarrollo del proyecto}
Durante las primeras semanas se determinó el entorno de desarrollo del proyecto, eligiendo las distintas herramientas que se iban a utilizar para el desarrollo del proyecto.

Una vez determinado todo el entorno se comenzó con el desarrollo, empezando con la creación del entorno virtual y la creación y conexión con la base de datos.

\subsection{Creación del entorno virtual}
Lo primero fue crear el entorno virtual con Flask en Visual Studio Code. Para ello, se creó una carpeta llamada 'env'
En el terminal de Visual Studio Code, una vez ubicado en la raíz del proyecto, se
ejecutó el comando: python3 -m venv env

 para crear el entorno virtual llamado 'env'.

Después de crear el entorno virtual al estar usando Windows se debía activar con el comando:
source env/bin/activate

Una vez se tenía el entorno virtual activado se instalaron las dependencias de flask en el entorno:
pip install Flask

De esta forma ya se podía iniciar la aplicación Flask con el comando: flask run.

\subsection{Creación y conexión de la base de datos}
La creación y conexión de la base de datos se hizo a través de la terminal de Visual Studio Code.
Para ello, lo primero fue:

\subsection{Creación del frontend}
Para crear el frontend, lo primero que se hizo fue desarrollar un prototipo de las interfaces de la aplicación, con el objetivo de diseñar una primera versión de cómo serían las interfaces. Este primer diseño, como se trataba de un prototipo, constó de interfaces sencillas e intuitivas para facilitar la comprensión y la navegación del usuario.

El desarrollo de las interfaces de la aplicación fue evolutivo. Para su diseño visual se utilizó CSS, ya que permitía personalizar de forma manual los templates HTML de forma sencilla y eficiente.

En un primer momento, se implementaron las interfaces de inicio de sesión y registro de nuevas cuentas de usuarios. Posteriormente, se diseñó la interfaz principal del menú de juegos para el usuario, que incluía una barra de búsqueda y filtros, así como tarjetas con la información de cada juego. Además, en la barra de navegación superior había un botón para solicitar el rol de profesor. Dependiendo del rol del usuario, la interfaz del menú de juegos cambiaba para mostrar funcionalidades específicas.

La interfaz para el menú de juegos del profesor, además de la barra de búsqueda y filtros, y las tarjetas con información de cada juego, incluía tres botones adicionales para añadir nuevos juegos, modificar juegos existentes y agregar archivos a cada juego. Esta interfaz se diferenciaba de la del usuario y la del administrador.

Por su parte, la interfaz del administrador también incluía un botón para la administración,lo que la diferenciaba de la del menú de juegos del usuario. En esta interfaz, el administrador podía realizar diferentes tareas de gestión.

Después, se trabajó en el desarrollo de interfaces específicas para la visualización y modificación de cada juego, con el fin de proporcionar una experiencia más detallada y personalizada para los usuarios.

Por último, se crearon interfaces para los usuarios con rol de administrador que les permitían gestionar usuarios, juegos y solicitudes en general. Además, se incluyeron las interfaces vara la valoración de los juegos por parte de los usuarios y una interfaz adicional para proporcionar información general sobre la aplicación y sus funcionalidades.

Una vez que se terminaron todos los templates junto con sus diseños, surgió un inconveniente con los estilos diseñados. El problema radicaba en que estos estilos no eran responsive y solo se veían correctamente cuando la pantalla estaba maximizada. Si el usuario reducía la ventana del navegador, todos los elementos se sobreponían unos encima de otros, lo que impedía la correcta visualización y usabilidad de la aplicación. Esto también significaba que la aplicación no era adecuada para su ejecución en distintos tipos de dispositivos.

Por ello, se empleó la librería de Bootstrap para lograr que la aplicación fuera responsive y que todos los elementos se visualizaran correctamente, ya fuera reduciendo la ventana del navegador, haciendo zoom o en otros tipos de dispositivos.
    
    (FOTO DE VISUALIZACION EN DISTINTOS DISPOSITIVOS)
    
\subsection{Creación del backend}
Para crear el backend, lo primero que se hizo fue establecer la conexión con la base de datos y crear las tablas necesarias, lo cual se llevó a cabo en el archivo database.py. 

Se crearon cuatro tablas que forman parte de la base de datos: 
\begin{enumerate}
    \item \textbf{Tabla de usuarios:} almacena los datos de los usuarios registrados en la aplicación, incluyendo su nombre de usuario, nombre, apellido, institución, contraseña y rol.
    \item \textbf{Tabla de juegos:} almacena los datos de los juegos registrados en la aplicación, incluyendo su nombre, descripción, idioma, enlace, puntuación, disciplina, naturaleza, precio, instrucciones para el jugador, instrucciones para el instructor, objetivo, espacio de control, objetivos principales y secundarios, estructura de sesiones, aspectos adicionales, la valoración del docente respecto al entretenimiento, aprendizaje, complejidad para el alumno, complejidad para el instructor, la URL del vídeo tutorial, los datos del usuario que añade o modifica el juego y los archivos de instrucciones y juego.
    \item \textbf{Tabla de solicitudes:} almacena los datos de las solicitudes para obtener el rol de profesor, incluyendo el estado de la solicitud, la fecha y el usuario que la realizó.
    \item \textbf{Tabla de valoraciones:} almacena los datos de las valoraciones que realizan los usuarios de los juegos, incluyendo la puntuación, el comentario, la fecha y los datos del usuario y juego al que se refiere la valoración.
\end{enumerate}

Después, se comenzó a desarrollar en el archivo app.py los diferentes métodos GET y POST para manejar las solicitudes y respuestas, así como los archivos juego.py y usuario.py para manejar las operaciones de inserción, eliminación y actualización de datos en la base de datos.

Para el desarrollo del login se utilizaron varios módulos de la librería Flask-Login, ya que esto facilitó toda la implementación. La clase LoginManager permitió gestionar la sesión del usuario en la aplicación, permitiendo a los usuarios iniciar sesión y mantenerse autenticados. Para ello, se creó una instancia de la clase LoginManager en la aplicación, lo que permitió usar sus métodos y atributos para configurar la autenticación y el inicio de sesión.

Además, en el caso de que un usuario intentara acceder a una ruta protegida sin estar autenticado, Flask redireccionaría al usuario a la vista de inicio de sesión, gracias a que se definió la propiedad login-view de la instancia de LoginManager creada anteriormente. 

        (FOTO DE lineas 54,55,56 -> App-py)

Con el módulo check-password-hash de la librería Werkzeug se permitía comprobar si la contraseña introducida para el inicio de sesión era correcta. En caso de que el usuario introdujera la contraseña de forma incorrecta durante 3 intentos consecutivos, se bloqueaba su cuenta durante 60 segundos para evitar ataques de fuerza bruta o intentos de acceso no autorizados. Este mecanismo de seguridad se implementó mediante la configuración de dos parámetros en el objeto app.config de Flask: MAX-LOGIN-ATTEMPTS, que establece el número máximo de intentos de inicio de sesión permitidos, y LOGIN-BLOCK-DURATION, que define la duración del bloqueo de la cuenta en segundos.

\subsection{Internacionalización}
Para implementar la internacionalización en mi aplicación, consideré dos opciones: el uso de la librería Flask-Babel o la inclusión de todas las traducciones en distintos archivos JSON.

Finalmente, opté por la segunda opción ya que me resultó más sencilla de implementar y gestionar debido a que solo necesitaba dos opciones de idioma. Flask-Babel ofrece características avanzadas, como la detección automática del idioma y de los formatos de fecha y hora, la funcionalidad de pluralización para la correcta traducción de los textos, y la capacidad de generar formularios de selección de idioma de forma automática, entre otras, por lo que consideré que era una solución más adecuada para proyectos más complejos.

Para cada template HTML se creó un archivo JSON, lo que permitió mantener todas las traducciones de cada página en un solo archivo, facilitando su llamado en el código.

    (FOTO DE LOS ARCHIVOS)

La aplicación ofrecía la posibilidad de visualizar el contenido en español o inglés, y se proporcionaron botones para que los usuarios seleccionen el idioma en la página de inicio. Si el usuario no seleccionaba un idioma, se utilizaba el español por defecto.

    (FOTO DE LOS BOTONES DE IDIOMAS)

\subsection{Despliegue con }

\section{Resolución de problemas}
Durante el desarrollo del proyecto, surgieron varios problemas que fueron resueltos con éxito.

\subsection{Problemas con ZenHub}
El primer problema que surgió fue con la plataforma de gestión de proyectos ZenHub. Al principio, utilicé esta herramienta para la gestión del proyecto, pero después de un mes recibí una notificación informándome de que la versión de prueba gratuita había finalizado.

Busqué alternativas que proporcionasen funcionalidades similares a las de ZenHub y finalmente opté por Zube, un sistema de gestión de proyectos gratuito integrado en GitHub que ofrece la planificación de sprints, la creación de tareas, la estimación de story points y la generación de gráficos e informes estadísticos de manera similar a ZenHub.

\subsection{Problemas con la máquina virtual}
\capitulo{6}{Trabajos relacionados}

En este apartado se presentan varios proyectos que guardan relación con el presente.

Aunque el objetivo específico de este proyecto ha dificultado la identificación de proyectos con el mismo propósito, se han encontrado TFGs de otras universidades que emplean herramientas similares a las que se utilizan en este.

Además, los tutores me proporcionaron un repositorio de juegos para el aula que resultaba muy similar al objetivo final de este proyecto.

Estos trabajos relacionados han facilitado la orientación del desarrollo y la aplicación de distintas herramientas y técnicas en este.

\section{The Gamification Repository}
\cite{Repositorio} Se trata de una aplicación para la gamificación en "Dirección de Proyectos” desarrollado por Lorenzo Borreguero Cortón y dirigido por Adolfo López-Paredes, de la Universidad de Valladolid.

Este repositorio recopila juegos en el aula relevantes para la docencia de materias de áreas de la Ingeniería de Organización y Organización de Empresas. 

La aplicación web \cite{wordpress} proporciona al usuario tres secciones diferentes:
\begin{itemize}
    \item\textbf{Gamificación:} proporciona al usuario información sobre el concepto de gamificación.
    \item\textbf{Buscador:} permite al usuario buscar los juegos de interés que se encuentran en la base de datos. Estos juegos se pueden buscar por el nombre, la asignatura relacionada o por el idioma. 
    
    Una vez buscado el juego por las opciones disponibles se muestran los distintos juegos que cumplen el criterio, proporcionando la url donde se obtiene el juego, su descripción, sus instrucciones y una descripción.
    \item\textbf{INSISOC Official Web:} redirecciona al usuario a la web oficial de INSISOC Grupo de Investigación Reconocido (GIR) de la Universidad de Valladolid, dedicado a la Ingeniería de los Sistemas Sociales.    
\end{itemize}

\section{Web corporativa e intranet para una consultora tecnológica}
\cite{ponce_web_2020} Se trata de un trabajo de fin de grado realizado por Pablo Sánchez Ponce en junio de 2020.

El objetivo de su proyecto consiste en el desarrollo de una aplicación web para la exposición de los servicios que proporciona una empresa, junto con la venta online de diversos cursos de formación para los clientes. 

Tanto mi proyecto como el trabajo de grado mencionado anteriormente tienen como objetivo la creación de una plataforma web que permita a los usuarios buscar, acceder y utilizar ciertos recursos. Ambos proyectos involucran la gestión de usuarios y recursos, con diferentes roles y permisos.

En su caso las tecnologías usadas para el desarrollo coinciden en el lenguaje de programación Python y el framework empleado Flask. Aunque, para el almacenamiento de datos se emplea el sistema MySQL a diferencia del mío en el que se emplea PostgreSQL.

\section{Desarrollo de Aplicación Web para Administración de Sensores Inalámbricos Mediante el Estándar BACnet}
\cite{pop_desarrollo_2020} Se trata de un trabajo de fin de grado realizado por Ciprian Ilut Pop en septiembre de 2020.

El objetivo de su proyecto consiste en el desarrollo de una aplicación web para la gestión de sensores a través del estándar BACnet.

En cuanto a las similitudes en objetivos, ambos proyectos tienen como objetivo el desarrollo de una aplicación web, y utilizan tecnologías similares, como Flask y PostgreSQL. Además, ambos proyectos tienen como objetivo proporcionar información y funcionalidades a los usuarios. 

Aunque el objetivo final de este proyecto es diferente al objetivo del propio, varias de las herramientas y técnicas que se emplean son las mismas. En este caso, para el almacenamiento de datos también se emplea el sistema de PostgreSQL. Para la implementación de la aplicación web se ha usado Flask como framework y Python como lenguaje de programación.

\section{Desarrollo de una aplicación Web RESTfull basada en Spring y Angular}
\cite{RESTfull_Jorge} Se trata de un trabajo de fin de grado realizado por Jorge Zapata Alburquerque.

El objetivo de su proyecto consiste en la creación de una aplicación web para la administración de currículums en una empresa. Su objetivo principal es mantener un registro actualizado de las habilidades de los empleados.

Los objetivos y funcionalidades son diferentes en cada proyecto, ya que mi aplicación se enfoca en la búsqueda y filtrado de juegos docentes, permitiendo a los usuarios puntuar los juegos, y su objetivo principal es mantener un registro de las habilidades de los empleados, enfocándose en la gestión de datos de la empresa y no en la interacción de los usuarios con la aplicación.

En cuanto a las similitudes, ambos proyectos utilizan PostgreSQL para la gestión de la base de datos. También se utilizan tecnologías web para el desarrollo de la aplicación, aunque mi proyecto utiliza Python con Flask y el trabajo de fin de grado de Zapata se basa en Spring y Angular.

\capitulo{7}{Conclusiones y Líneas de trabajo futuras}

En este apartado, se expondrán las conclusiones obtenidas tras finalizar el desarrollo del proyecto, así como las posibles líneas de trabajo que podrían llevarse a cabo en el futuro.

\section{Conclusiones}
Una vez finalizado el proyecto, se puede considerar que se han cumplido satisfactoriamente los objetivos generales, técnicos y personales establecidos inicialmente. La aplicación desarrollada ha logrado ser una herramienta efectiva y accesible para el aprendizaje a través de juegos docentes, cumpliendo con los requisitos establecidos.

Además, se han añadido funcionalidades adicionales que enriquecen la experiencia de usuario y amplían las capacidades de la aplicación. Estas mejoras incluyen la implementación de filtros de búsqueda y un sistema de gestión completo para administrar juegos, usuarios y solicitudes.

Durante el desarrollo del proyecto, se han adquirido y fortalecido conocimientos en diversas herramientas y lenguajes. Se ha mejorado la comprensión y habilidades en HTML, CSS y se ha utilizado el framework Bootstrap para agilizar el desarrollo y lograr diseños responsive.

En cuanto a Flask, se ha profundizado en su uso y se ha aprovechado su versatilidad para el desarrollo de aplicaciones web. Se ha trabajado en la definición de rutas, diseño de vistas y utilización de plantillas para generar contenido dinámico. También se ha implementado formularios y validaciones para mejorar la interacción con el usuario.

En Python, se han ampliado los conocimientos y habilidades, especialmente en el contexto del desarrollo web. Se han utilizado diversas funcionalidades y bibliotecas para la manipulación de datos, interacción con la base de datos PostgreSQL y la implementación de la lógica de negocio de la aplicación.

Además, se ha adquirido experiencia en la metodología ágil Scrum y se han utilizado herramientas como Zube para la planificación y gestión eficiente de tareas. 

En resumen, el proyecto ha permitido ampliar y aplicar los conocimientos adquiridos durante el grado, desarrollando una aplicación web educativa con éxito.

\section{Líneas de trabajo futuras}
Algunas ideas para futuras líneas de desarrollo de la aplicación web pueden ser:

\begin{itemize}
    \item Incorporar una función de mensajería interna para que los usuarios puedan comunicarse entre sí, de forma que los docentes puedan compartir ideas y conocimientos sobre los distintos juegos docentes.
    \item Mejorar la experiencia de búsqueda de juegos. Agregando filtros adicionales como el tema del juego o su nivel de dificultad, de forma que los usuarios puedan encontrar juegos más específicos.
    \item Implementar una página personalizada para los profesores donde puedan ver sus juegos favoritos o los juegos que han creado.
    \item Añadir una sección de informes para que los administradores puedan ver distintas estadísticas sobre el uso de la plataforma. Incluyendo información como la cantidad de usuarios activos, la cantidad de juegos subidos en un período determinado de tiempo y las solicitudes de rol aceptadas o rechazadas.
    \item Integrar la plataforma con las redes sociales, permitiendo que los usuarios puedan compartir los juegos docentes.
    \item Incluir un foro de discusión para que los usuarios puedan hacer preguntas y compartir ideas sobre cómo utilizar los juegos educativos en el aula.
\end{itemize}


\bibliographystyle{plain}
\bibliography{bibliografia}

\end{document}
